\section{Related Work}
\label{sec:related}

\subsection{Pharmaceutical regulatory evidence and shortages}

The FDA Drug Shortages Task Force identified weak incentives for low-margin products, limited market recognition of mature quality systems, and logistical and regulatory barriers to recovery as recurring root causes \cite{fda2019shortages}. FDA shortage records use defined categories that include current good manufacturing practice requirements, regulatory delay, active- or inactive-ingredient shortage, discontinuation, shipping delay, and increased demand \cite{fdaShortageFAQ}. The openFDA interface publishes machine-readable shortage fields and harmonized identifiers, which support reproducible collection but do not turn database fields into regulatory conclusions \cite{openfdaShortages}. FDA's draft risk-management-plan guidance links proactive supply-risk identification to ICH Q9 principles, while explicitly retaining draft status \cite{fdaShortageRMP2022,ichQ9R1}.

Regulatory evidence is distributed across authorities, document families, and lifecycle stages. ICH Q10 treats knowledge management and quality risk management as enablers throughout the product lifecycle \cite{ichQ10}. FDA's quality-systems guidance explains how a modern quality system can operate consistently with current good manufacturing practice \cite{fdaQualitySystems2006}. Annex 11 adds computerized-system lifecycle controls \cite{ecAnnex11}. Our system preserves the authority, document, section, and source location of each passage because a retrieved statement must remain distinguishable from a generated summary or a graph relation.

Recent domain studies provide close precedents. Gray et al. classified free regulatory text into standard labeling sections \cite{gray2023regulatory}. Saraswat et al. evaluated question similarity for pharmaceutical regulatory search using expert-annotated entailment pairs \cite{saraswat2023regulatory}. Proestel et al. studied semantic search over FDA guidance, and Nanua et al. evaluated RAG against laboratory regulations \cite{proestel2025semantic,nanua2025raven}. These studies support the practical relevance of regulatory retrieval. Our evaluation differs by separating source-chunk ranking from structured shortage-path recovery and by testing whether added semantic and graph components improve retrieval on a fixed corpus.

\subsection{Lexical, dense, and neural retrieval}

BM25 derives a probabilistic lexical score from term frequency, inverse document frequency, and length normalization \cite{robertson2009bm25}. Dense passage retrieval instead learns query and passage encoders whose vectors can be searched efficiently \cite{karpukhin2020dpr}. The two approaches respond to different signals. A dense model can bridge paraphrases, while a lexical model can exploit exact guideline numbers, table labels, product terms, and defined phrases. Our experiments therefore treat BM25 as a full baseline rather than a preliminary component that semantic retrieval is assumed to surpass.

HyDE addresses zero-shot dense retrieval by generating a hypothetical document for a query and using its embedding to locate real passages \cite{gao2023hyde}. We adapt the same general idea on the index side: generated questions and summaries point back to source chunks and never count as evidence. Youtu-Embedding, the frozen dense encoder used in our experiments, was trained through a framework designed to balance information retrieval and semantic textual similarity objectives \cite{zhang2025codiemb}. We evaluate the resulting representations in this corpus rather than transferring benchmark claims to pharmaceutical text.

Neural rerankers apply richer query-document interaction after a first-stage retriever has reduced the candidate set. Sequence-to-sequence rerankers can interpret relevance-label token probabilities as ranking scores \cite{nogueira2020monot5}; multilingual long-context rerankers extend this pattern beyond short English passages \cite{zhang2024mgte}. Qwen3-Reranker provides generative rerankers at several parameter scales \cite{zhang2025qwen3embedding}. We use its 0.6B variant only in a locked supplementary analysis over BM25's top 50. This design tests candidate reordering without allowing the reranker to recover a source chunk that BM25 failed to retrieve.

\subsection{Hierarchical and graph-based retrieval}

RAPTOR recursively embeds, clusters, and summarizes text into a retrieval tree to expose multiple levels of abstraction \cite{sarthi2024raptor}. GraphRAG extracts entity graphs and community summaries to answer corpus-level questions \cite{edge2024graphrag}. Those objectives differ from exact regulatory passage retrieval. A community summary may help explain a corpus but cannot replace the source clause used to justify a regulated decision.

Knowledge graphs provide explicit entities, relations, identity choices, schemas, and provenance \cite{hogan2021knowledgegraphs}. We use these properties to represent shortage events and document structure. We do not assume that graph connectivity is equivalent to textual relevance. Reciprocal rank fusion (RRF) offers a score-scale-independent way to combine rankings \cite{cormack2009rrf}; our experiments include both unconditional RRF and constrained variants so that a failed fusion remains visible rather than being removed after inspection.
